\chapter{LIMITATIONS, CONCLUSION AND FUTURE WORK}\nonum

\section{Limitations}
\begin{itemize}
\item \textbf{No clipping.} A triangle with a corner behind the camera is skipped completely, not cut. In the demo scene this skipped 359 of 15,412 triangles when culling was off.
\item \textbf{Flat shading.} Each triangle has one colour, so curved objects show their flat faces.
\item \textbf{One-sided planes.} A plane has one front side and disappears when seen from behind with culling on.
\item \textbf{Visvalingam--Whyatt is $O(n^2)$.} It is fine for the lab but would be slow for millions of points.
\item \textbf{One run, one machine.} The benchmark was run once in headless Chrome. The relative results are useful, but the exact milliseconds will differ on other computers.
\item \textbf{Bresenham is not faster in JavaScript,} as the benchmark shows.
\item \textbf{The interface has no automated tests.} We checked it by hand, in Chrome on macOS only.
\item \textbf{Lighting, fog and bloom} are simple extras and are not among the four required algorithms.
\end{itemize}

\section{Conclusion}
RenderLab has the four required algorithms (Bresenham, back-face culling, alpha compositing and Visvalingam--Whyatt) as plain TypeScript with no graphics library. Each has a lab that shows the formulas, the numbers at every step, and a comparison. Other labs cover transformations, determinants, quaternions, the outer product, the sine rule, line intersection, bases and handedness. All 487 tests pass, and all 195 controls passed our click-through after two layout fixes. The benchmark gave some results we did not expect: DDA was faster than Bresenham in JavaScript, and Visvalingam--Whyatt was about 34 times slower than RDP at 2,000 points.

\begin{table}[H]
\centering
\caption{Contribution of each student}
\label{tab:contrib}
\small
\begin{tabularx}{\linewidth}{@{}p{3.2cm}L@{}}
\toprule
Name & Work\\
\midrule
Yasharth Singh & TODO: fill in before submission\\
Rakshit Suneja & TODO: fill in before submission\\
\bottomrule
\end{tabularx}
\end{table}

\section{Future Work}
\begin{itemize}
\item Port the core to Rust and WebAssembly. Bresenham's whole-number arithmetic should then show the speed benefit that JavaScript does not give, and the benchmark can be repeated to check.
\item Use a heap in Visvalingam--Whyatt to make it $O(n\log n)$.
\item Cut triangles at the near plane (clipping) so they do not disappear close to the camera.
\item Add smooth shading and textures.
\item Add automated tests for the interface.
\end{itemize}
